\begin{abstract}
Los canales ausentes en registros de EEG representan un desafío significativo en neurofisiología clínica y en pipelines de Interfaces Cerebro-Computadora, ya que los errores de interpolación pueden distorsionar la morfología, latencia y contenido espectral de biomarcadores posteriores. Este artículo presenta un marco riguroso y reproducible para la evaluación comparativa de reconstrucción de canales EEG bajo condiciones de canales faltantes. Utilizando un pipeline unificado de optimización y validación en tres conjuntos de datos (PhysioNet, BCI Competition IV 2a y MNE Sample), comparamos métodos basados en grafos con regularización temporal (TRSS, Tikhonov) contra líneas base geométricas exhaustivamente optimizadas (Spherical Spline, RBF). En lugar de proponer un algoritmo fundamentalmente novedoso, nuestra contribución principal es metodológica: demostramos que la optimización rigurosa de hiperparámetros mediante Optuna cierra sustancialmente las brechas de rendimiento entre enfoques clásicos y modernos. Bajo esta evaluación justa, los métodos de regularización temporal mantienen márgenes competitivos y estadísticamente significativos. TRSS demuestra fortaleza particular en la preservación de dinámicas fisiológicas (Potenciales Relacionados con Eventos, Densidad Espectral de Potencia) bajo pérdida severa de canales contiguos, estableciendo un punto de referencia para el rendimiento realista de interpolación EEG en todas las clases de métodos.
\end{abstract}

\begin{IEEEkeywords}
reconstrucción EEG, procesamiento de señales de grafos, interpolación, Optuna, Potenciales Relacionados con Eventos, regularización temporal
\end{IEEEkeywords}
